\pageoriginale

\clearpage
\SGAExposeBookmark{XVI}{La fórmula de Milnor}
\begin{center}
{\large SGA 7 --- EXPOSÉ XVI}\par
\vspace{2em}
{\Large\underline{La fórmula de Milnor}}\par
\vspace{1.5em}
{\large por P. Deligne}
\end{center}

\medskip
En esta exposición se transpone a característica $p$ la fórmula dada por
Milnor [3] (véase también [2]), que permite en ciertos casos calcular
las dimensiones de espacios de ciclos evanescentes. La demostración
es de naturaleza global y solo permite tratar el caso ``geométrico''
de una base que sea un trait henseliano de igual característica con
cuerpo residual perfecto.

\section*{1. Enunciado del problema}

\subsubsection*{1.1.}
Sean $S$ un esquema regular puramente de dimensión uno, $s$ un punto
cerrado de $S$ \underline{con cuerpo residual perfecto},
$f:X\to S$ un morfismo plano de tipo finito, $x_0$ un punto cerrado
de la fibra $X_s$ y $\bar x_0$ un punto geométrico de $X$ localizado
en $x_0$. Supongamos que
\begin{enumerate}
\item[(a)] $x_0$ es un \underline{punto aislado de no lisitud} de
$f$: existen un entero $n$ y un entorno $U$ de $x_0$, puramente de
dimensión relativa $n$, tal que $U-\{x_0\}$ es liso sobre $S$;
\item[(b)] $X$ es \underline{regular} en $x_0$.
\end{enumerate}

Bajo estas hipótesis, $X$ es una intersección completa relativa sobre
$S$ en $x_0$. Remitimos a VI, 3.10 para la definición de
$D^1_{X/S}$. En el entorno $U$ de $x_0$ se tiene
\[
D^1_{X/S}
=\mathcal H^1\!\left(R\mathcal Hom(L\Omega_{X/S},\mathcal O_X)\right)
=\mathcal Ext^1(\Omega^1_{X/S},\mathcal O_X);
\]
en este entorno, $D^1_{X/S}$ tiene longitud finita y soporte contenido
en $\{x_0\}$.

\pageoriginale

\subsubsection*{Definición 1.2.}
\underline{El número de Milnor
$\mu=\mu(X/S,\bar x_0)$ de $X/S$ en $\bar x_0$ es la longitud del
$\mathcal O_{X,x_0}$-módulo definido por $D^1_{X/S}$.}

\subsubsection*{Ejemplo 1.3.}
Tomemos por $S$ una curva lisa sobre un cuerpo algebraicamente cerrado
$k$. Reemplazando $X$ por un entorno de $x_0$, podemos suponer que
$X$ es liso sobre $k$. El complejo cotangente se reduce entonces a la
resolución siguiente de $\Omega^1_{X/S}$:
\[
L\Omega_{X/S}
=\bigl[f^*\Omega^1_{S/k}\longrightarrow\Omega^1_{X/k}\bigr].
\]
Elijamos una coordenada local $t:S\to\mathbb A^1_k$ sobre $S$ en
$s$ (con $t$ étale en $s$), designemos de nuevo por $f$ el compuesto
$t\circ f$, y sea $(x_i)_{1\leq i\leq n+1}$ un sistema de
coordenadas locales sobre $X$ en $x_0$. Entonces, en $x_0$, el módulo
$f^*\Omega^1_{S/k}$ tiene base $df$, mientras que
$\Omega^1_{X/k}$ tiene por base los $dx_i$. Por consiguiente,
$D^1_{X/S}$ se identifica con el cociente de $\mathcal O_X$ por el
ideal engendrado por los $\partial f/\partial x_i$, y
\begin{equation}\tag{1.3.1}
\mu=\dim_k\!\left(
\mathcal O_{X,x_0}/
\left(\frac{\partial f}{\partial x_1},\ldots,
      \frac{\partial f}{\partial x_{n+1}}\right)\right).
\end{equation}
Esta es la definición que figura en Milnor.

\subsubsection*{Ejemplo 1.4.}
Supongamos que $X$ es el subesquema cerrado de $\mathbb A_S^{n+1}$
definido por una ecuación $g$, y sea
$i:X\hookrightarrow\mathbb A_S^{n+1}$ la inclusión. El complejo
cotangente se reduce entonces a
\[
\bigl[(g)/(g)^2\longrightarrow
i^*\Omega^1_{\mathbb A^{n+1}/S}\bigr],
\]
y $D^1_{X/S}$ es isomorfo al cociente de $\mathcal O_X$ por el ideal
engendrado por $(\partial g/\partial x_i)_{1\leq i\leq n+1}$; o,
equivalentemente, al cociente de $\mathcal O_{\mathbb A^{n+1}}$ por
el ideal engendrado por $g$ y los $\partial g/\partial x_i$.

\subsubsection*{Ejemplo 1.5.}
Tomemos por $S$ el espectro de un anillo de valoración discreta $V$
con cuerpo residual perfecto y cuerpo de fracciones $K$, y sea $X$
la normalización de $S$

\pageoriginale

en una extensión separable $K'$ de $K$. Por 1.4 y [5], III, 7.14,
p.~68, $D^1_{X/S}$ es entonces isomorfo al cociente de
$\mathcal O_X$ por la diferente $\mathfrak d_{X/S}$. Si $\bar x_0$
es un punto geométrico de $X_s$, localizado en $x_0\in X_s$, su
número de Milnor es, por tanto,
\begin{equation}\tag{1.5.1}
\mu=\operatorname{longitud}
\bigl(\mathcal O_{X,x_0}/\mathfrak d_{X/S}\bigr).
\end{equation}
Sea $\bar s$ un punto geométrico de $S$ localizado en el punto cerrado
$s$, y sea $d_{K'/K}\in V$ un generador del ideal discriminante, de
valoración $v(d_{K'/K})$. De (1.5.1) se deduce que
\begin{equation}\tag{1.5.2}
v(d_{K'/K})=\sum_{x\in X(\bar s)}\mu(X/S,x).
\end{equation}

\subsubsection*{1.6.}
Conservemos las hipótesis de 1.1. Sea $\bar s$ el punto geométrico de
$S$ imagen de $\bar x_0$, sea $S_{(\bar s)}$ el espectro del
henselizado estricto de $\mathcal O_{S,s}$ en $\bar s$, y sea
$X_{(\bar s)}$ la imagen inversa de $X$ sobre $S_{(\bar s)}$.
Sea $\bar\eta$ el espectro de una clausura algebraica del cuerpo de
fracciones $k(\eta)$ de $S_{(\bar s)}$, y pongamos
\[
I=\operatorname{Gal}(k(\bar\eta)/k(\eta)).
\]

Sea $\ell$ un número primo distinto de la característica residual de
$S$. En un entorno de $\bar x_0$ en $X_{\bar s}$, los haces de ciclos
evanescentes
\[
\Phi^i_{\bar\eta}(\mathbb Z/\ell)
=R^i\Phi_{\bar\eta}(\mathbb Z/\ell)
\]
relativos al haz constante $\mathbb Z/\ell$ sobre $X_{(\bar s)}$
están concentrados en $\bar x_0$. Además, el grupo de inercia $I$
actúa sobre los espacios vectoriales sobre $\mathbb Z/\ell$
\[
\Phi^i_{\bar\eta}(\mathbb Z/\ell)_{\bar x_0}.
\]
Estos espacios son de dimensión finita (XIII, 2.4.2).

\subsubsection*{1.7.}
Sea $I_1$ un cociente finito del grupo de inercia. Recordemos que el
\underline{carácter de Swan} $\operatorname{Sw}$ de $I_1$ es la
diferencia $a-u$ entre el carácter de Artin $a$ y el carácter $u$ de
la representación de aumentación. Recordemos también que, salvo
isomorfismo, existe un único $\mathbb Z_\ell[I_1]$-módulo proyectivo
$S(I_1)$ de carácter $\operatorname{Sw}$ (véase [6]).

\pageoriginale

Sea $V$ un espacio vectorial de dimensión finita sobre
$\mathbb Z/\ell$, provisto de una acción
$\rho:I\to\operatorname{GL}(V)$, y supongamos que $\rho$ se factoriza
por $I_1$. La \underline{dimensión salvaje} de $V$ es entonces
\begin{equation}\tag{1.7.1}
\dim\operatorname{Sw}(V)
=\dim_{\mathbb Z/\ell}
\operatorname{Hom}_{I_1}\!\left(
S(I_1)\otimes_{\mathbb Z_\ell}
\mathbb Z/\ell,V\right).
\end{equation}
Este entero es independiente de la elección de $I_1$, pues para un
subgrupo normal $H$ de $I_1$ se tiene $S(I_1/H)\sim S(I_1)_H$.

Ponemos además
\[
\dim\operatorname{tot}(V)
=\dim\operatorname{Sw}(V)+\dim_{\mathbb Z/\ell}(V)
\]
y llamamos a este número la \underline{dimensión total} de $V$.

\subsubsection*{Definición 1.8.}
\underline{Con la terminología de 1.6 y 1.7, se llama número de ciclos
evanescentes de $X/S$ en $\bar x_0$ (respectivamente, número salvaje
de ciclos evanescentes; respectivamente, número total de ciclos
evanescentes) al entero}
\[
(-1)^n\dim\Phi_{\bar x_0}
\underset{\mathrm{def}}{=}
(-1)^n\sum_i(-1)^i
\dim\Phi^i_{\bar\eta}(\mathbb Z/\ell)_{\bar x_0},
\]
\[
(-1)^n\dim\operatorname{Sw}\Phi_{\bar x_0}
\underset{\mathrm{def}}{=}
(-1)^n\sum_i(-1)^i
\dim\operatorname{Sw}
\Phi^i_{\bar\eta}(\mathbb Z/\ell)_{\bar x_0},
\]
\[
\begin{aligned}
(-1)^n\dim\operatorname{tot}\Phi_{\bar x_0}
&\underset{\mathrm{def}}{=}
(-1)^n\sum_i(-1)^i
\dim\operatorname{tot}
\Phi^i_{\bar\eta}(\mathbb Z/\ell)_{\bar x_0}\\
&=(-1)^n\bigl(
\dim\Phi_{\bar x_0}+\dim\operatorname{Sw}\Phi_{\bar x_0}
\bigr).
\end{aligned}
\]

\subsubsection*{Conjetura 1.9.}
\underline{Bajo las hipótesis de 1.1, el número de Milnor $\mu$ de
$X/S$ en $\bar x_0$ es igual al número total de ciclos evanescentes
de $X/S$ en $\bar x_0$.}

\subsubsection*{Observación 1.10.}
Sea $S$ un trait henseliano con cuerpo residual imperfecto y sea
$f:X\to S$

\pageoriginale

un morfismo liso salvo en $x_0\in X_s$. Supongamos que $X$ es regular
en $x_0$. La definición 1.2 del número de Milnor $\mu$ sigue teniendo
sentido bajo estas hipótesis, pero ya no tengo una interpretación
geométrica de $\mu$ que proponer.

\subsubsection*{1.11.}
En el resto de esta sección, supongamos que
$S=\operatorname{Spec}(V)$ es un trait henseliano con cuerpo residual
\underline{algebraicamente cerrado} y punto cerrado $s$, que
$f:X\to S$ es un morfismo plano de tipo finito, puramente de dimensión
relativa $n$, que $X$ es regular, que $x_0$ es un punto cerrado de
$X_s$ y que $X-\{x_0\}$ es liso sobre $S$.

Esto no restringe la generalidad.

\subsubsection*{Proposición 1.12.}
\underline{Si $n=0$ (es decir, si $X/S$ es finito), la Conjetura 1.9
es verdadera para $(X/S,x_0)$.}

Podemos suponer que $X$ es la normalización de $S$ en una extensión
finita $K'$ del cuerpo de fracciones $K$. Por 1.5,
$\mu=v(d_{K'/K})$.

La extensión $K'$ está definida por un subgrupo abierto $H$ del grupo
de inercia $I$. Sea $H'\subset H$ un subgrupo abierto normal de $I$,
y pongamos $I'=I/H'$. Sea $s_{I/H}$ el carácter de la representación
de permutación de $I'$ sobre $I/H$, sea $a_{I'}$ el carácter de Artin,
sea $\operatorname{Sw}_{I'}$ el carácter de Swan y sea $u_{I'}$ el
carácter de la representación de aumentación. Por la
``Führerdiskriminantenproduktformel'' de Artin, el ideal discriminante
es el conductor de Artin de $s_{I/H}$ ([5], VI, 3, Corolario 1,
p.~111):
\[
\mu=(a_{I'},s_{I/H})
=|I'|^{-1}\sum_{i\in I'}a_{I'}(i)s_{I/H}(i).
\]

\pageoriginale

Se tiene
\[
a_{I'}=u_{I'}+\operatorname{Sw}_{I'},
\qquad
(u_{I'},s_{I/H})=\#(I/H)-1,
\]
y
\begin{align*}
(\operatorname{Sw}_{I'},s_{I/H})
&=\dim_{\mathbb Z/\ell}\operatorname{Hom}_I\!\left(
S(I')\otimes_{\mathbb Z_\ell}\mathbb Z/\ell,
(\mathbb Z/\ell)^{I/H}\right)\\
&=\dim_{\mathbb Z/\ell}\operatorname{Hom}_I\!\left(
S(I')\otimes_{\mathbb Z_\ell}\mathbb Z/\ell,
\operatorname{Coker}\bigl(
\mathbb Z/\ell\longrightarrow(\mathbb Z/\ell)^{I/H}\bigr)
\right).
\end{align*}
Por tanto,
\[
\mu=\#(I/H)-1+
\dim\operatorname{Sw}\operatorname{Coker}\bigl(
\mathbb Z/\ell\longrightarrow(\mathbb Z/\ell)^{I/H}\bigr).
\]
Se comprueba inmediatamente que
$\Phi^i_{x_0}=0$ para $i\ne0$ y que
\[
\Phi^0_{x_0}=\operatorname{Coker}\bigl(
\mathbb Z/\ell\longrightarrow(\mathbb Z/\ell)^{I/H}\bigr);
\]
la fórmula se deduce de ello.

\subsubsection*{Proposición 1.13.}
\underline{Si $X_s$ presenta en $x_0$ una singularidad cuadrática
ordinaria, la Conjetura 1.9 es verdadera para $(X/S,x_0)$.}

\medskip
\noindent\underline{Caso 1: $n=\dim(X_s)=2m-1$.}

El problema es local para la topología étale sobre $X$. Podemos, por
tanto, suponer que $\sisi{X_s}{X}$ es el subesquema cerrado de
$\mathbb A_S^{n+1}$ definido por una ecuación
\[
f=\sum_{i=1}^{m}x_i x_{i+m}+a=0,
\]
donde $a$ pertenece al ideal maximal de $V$ y $x_0$ es el origen de
las coordenadas. Además, $X$ es regular en $x_0$ si y solo si $a$ es
un uniformizante. Por 1.4,
\[
\mu=\operatorname{longitud}\!\left(
\mathcal O_{\mathbb A_S^{2m},x_0}/
\left(f,\frac{\partial f}{\partial x_i}\right)\right)=1.
\]

\pageoriginale

Por el cálculo de los ciclos evanescentes cuadráticos ordinarios del
Exposé XV,\footnote{Después de la palabra impresa ``Por'', la fuente
deja en blanco la referencia esperada. La concentración y la acción de
inercia enunciadas son el cálculo de ciclos evanescentes cuadráticos
ordinarios establecido en el Exposé XV.}
$\Phi^i_{x_0}=0$ para $i\ne n$, y $\Phi^n_{x_0}$ tiene dimensión uno,
con acción trivial de la inercia. El número salvaje de ciclos
evanescentes es, por tanto, cero, y el número total es $1$, lo que
verifica 1.9.

\medskip
\noindent\underline{Caso 2: $n=\dim(X_s)=2m$ y
$\operatorname{car}(k(s))\ne2$.}

Bajo estas hipótesis podemos tomar $X\subset\mathbb A_S^{n+1}$
definido por la ecuación
\[
f=\sum_{i=1}^{m}x_i x_{i+m}+x_{2m+1}^2+a=0,
\]
donde $a$ es un uniformizante. De nuevo $\mu=1$,
$\Phi^i_{x_0}=0$ para $i\ne n$, y
$\sisi{\Phi^i_{x_0}}{\Phi^n_{x_0}}$ tiene dimensión uno. Esta vez el
grupo de inercia actúa por $\pm1$ sobre $\Phi^n_{x_0}$, de modo
moderado. El número salvaje de ciclos evanescentes vuelve a ser cero,
y 1.9 queda verificada.

\medskip
\noindent\underline{Caso 3: $n=\dim(X_s)=2m$ y
$\operatorname{car}(k(s))=2$.}

Esta vez podemos tomar $X\subset\mathbb A_S^{n+1}$ definido por
\[
f=\sum_{i=1}^{m}x_i x_{i+m}
 +(x_{2m+1}^2+a x_{2m+1}+b)=0,
\]
donde $a$ y $b$ pertenecen al ideal maximal. El esquema $X$ es regular
si y solo si $b$ es un uniformizante, y su fibra genérica es lisa si
y solo si $a^2-4b\ne0$.

El número de Milnor de $X/S$ coincide con el de
$X_0\subset\mathbb A_S^1$ definido por
\begin{equation}\tag{1}
x^2+a x+b=0.
\end{equation}
Asimismo, $\Phi^i_{x_0}=0$ para $i\ne n$, $\Phi^n_{x_0}$ tiene
dimensión uno e $I$ actúa sobre él por $\pm1$; el subgrupo de índice
dos de $I$ que actúa trivialmente corresponde a la extensión de $K$
definida por (1). El número total de ciclos evanescentes de

\pageoriginale

$X/S$ en $x_0$ coincide, por tanto, con el relativo a $X_0/S$, y la
validez de 1.9 en este caso se deduce de 1.12. El número de Milnor
viene dado por (1.5):
\begin{equation}\tag{1.13.1}
\mu=v(a^2-4b).
\end{equation}

\section*{2. El caso ``geométrico''}

Nos proponemos demostrar que la Conjetura 1.9 es verdadera en igual
característica. El resultado geométrico clave es el siguiente. La
posibilidad de utilizarlo me fue sugerida por conversaciones con
B. Saint-Donat.

\subsubsection*{Proposición 2.1.}
\underline{Sean $S$ una curva proyectiva lisa y conexa sobre un cuerpo
algebraicamente cerrado $k$, $X$ una variedad propia y lisa sobre
$k$, puramente de dimensión $n+1$, y $f:X\to S$ un morfismo liso
fuera de un conjunto finito $F$ de puntos de $X$. Entonces las
características de Euler--Poincaré, en cohomología étale, de $S$, de
$X$ y de la fibra genérica geométrica $X_{\bar\eta}$ de $f$
satisfacen}
\begin{equation}\tag{2.1.1}
\chi(X)=\chi(S)\chi(X_{\bar\eta})
-(-1)^n\sum_{x\in F}\mu(X/S,x).
\end{equation}

Se sabe que la característica de Euler--Poincaré de $X$ coincide con
la traza (o integral sobre $X$) de la última clase de Chern del fibrado
tangente:
\[
\chi(X)=\int_Xc_{n+1}(\Theta_X)
=(-1)^{n+1}\int_Xc_{n+1}(\Omega_X^1).
\]
Pongamos
\[
\Omega_X^1=\left(\Omega_X^1\otimes
(f^*\Omega_S^1)^{\otimes-1}\right)\otimes f^*\Omega_S^1.
\]
La clase total de Chern del haz inversible $f^*\Omega_S^1$ es la
imagen inversa de la clase total de Chern de $\Omega_S^1$ sobre $S$:

\pageoriginale

\[
c^*(f^*\Omega_S^1)=1+f^*c^1(S,\Omega_S^1),
\]
y las potencias respecto del producto cup de $c^1$ son nulas:
\[
\bigl(f^*c^1(S,\Omega_S^1)\bigr)^n
=f^*\bigl(c^1(S,\Omega_S^1)^n\bigr)=0
\qquad(n>1).
\]
De ello se deduce que
\begin{align*}
c^{n+1}(\Omega_X^1)
={}&c^n\!\left(\Omega_X^1\otimes
(f^*\Omega_S^1)^{\otimes-1}\right)
\smile f^*c_1(S,\Omega_S^1)\\
&+c^{n+1}\!\left(\Omega_X^1\otimes
(f^*\Omega_S^1)^{\otimes-1}\right).
\end{align*}

Si $S$ tiene género $g$, entonces $c_1(S,\Omega_S^1)$ es
$(2g-2)$ veces la clase de cohomología de cualquier punto cerrado
$s$ de $S$. Por tanto, si $X_s$ es no singular,
\begin{align*}
\mathrm I
&=\int_Xc^n\!\left(\Omega_X^1\otimes
(f^*\Omega_S^1)^{\otimes-1}\right)
\smile f^*c_1(S,\Omega_S^1)\\
&=(2g-2)\int_Xc^n\!\left(\Omega_X^1\otimes
(f^*\Omega_S^1)^{\otimes-1}\right)
\smile f^*\operatorname{cl}(s)\\
&=(2g-2)\int_{X_s}c^n\!\left(
X_s,\Omega_X^1\otimes(f^*\Omega_S^1)^{\otimes-1}
\otimes\mathcal O_{X_s}\right).
\end{align*}
Además, sobre $X_s$ el haz $f^*\Omega_S^1$ es isomorfo a
$\mathcal O_{X_s}$, y
$\Omega_X^1\otimes_{\mathcal O_X}\mathcal O_{X_s}$ figura en una
sucesión exacta
\[
0\longrightarrow f^*\Omega_S^1\otimes\mathcal O_{X_s}
\longrightarrow\Omega_X^1\otimes\mathcal O_{X_s}
\longrightarrow\Omega_{X_s}^1\longrightarrow0.
\]
Por consiguiente,
\[
c^n\!\left(X_s,\Omega_X^1\otimes
(f^*\Omega_S^1)^{\otimes-1}\otimes\mathcal O_{X_s}\right)
=c^n(X_s,\Omega_{X_s}^1),
\]
y
\[
\mathrm I=(2g-2)\int_{X_s}c^n(X_s,\Omega_{X_s}^1)
=(-1)^{n+1}\chi(S)\chi(X_s).
\]

\pageoriginale

La aplicación canónica
\[
f^*:f^*\Omega_S^1\longrightarrow\Omega_X^1
\]
se identifica con una sección $a$ de
$\Omega_X^1\otimes(f^*\Omega_S^1)^{\otimes-1}$. Una vez elegidos
sistemas de coordenadas locales sobre $S$ y $X$, lo que identifica
este fibrado con $\mathcal O_X^{n+1}$, las coordenadas de $a$ son las
derivadas parciales $\partial f/\partial x_i$. Por tanto, $a$ solo se
anula sobre $F$, y por 1.3 la multiplicidad de su cero en $x\in F$ es
igual al número de Milnor $\mu(X/S,x)$.

Recordemos que si una sección $a$ de un fibrado vectorial $G$ de rango
$d$ sobre un $k$-esquema propio y liso $X$, puramente de dimensión
$d$, solo se anula en puntos aislados, entonces
$\int_Xc^d(G)$ es igual al número de ceros de $a$, contados con sus
multiplicidades. Aquí se tiene
\[
\mathrm{II}=\int_Xc^{n+1}\!\left(
\Omega_X^1\otimes(f^*\Omega_S^1)^{\otimes-1}\right)
=\sum_{x\in F}\mu(X/S,x).
\]
Reuniendo los términos, se obtiene
\[
\chi(X)=(-1)^{n+1}(\mathrm I+\mathrm{II})
=\chi(S)\chi(X_s)-(-1)^n\sum_{x\in F}\mu(X/S,x),
\]
de donde se deduce 2.1, pues
$\chi(X_s)=\chi(X_{\bar\eta})$.

\subsubsection*{Corolario 2.2.}
\underline{Bajo las hipótesis precedentes, si $n(X/S,x)$ es el número
total de ciclos evanescentes de $X/S$ en $x$ (1.8), entonces}
\[
\sum_{x\in F}\mu(X/S,x)
=\sum_{x\in F}n(X/S,x).
\]

La sucesión espectral de Leray para $f$ proporciona
\begin{align*}
\chi(X)
&=\sum_{p,q}(-1)^{p+q}
\dim_{\mathbb Z/\ell}H^p(S,R^qf_*\mathbb Z/\ell)\\
&=\sum_q(-1)^q\chi(S,R^qf_*\mathbb Z/\ell).
\end{align*}

\pageoriginale

Para todo haz constructible $G$ de $\mathbb Z/\ell$-módulos sobre
$S$ se tiene ([4] o SGA 5, X, 7.2)
\[
\chi(S,G)=\chi(S)\dim_{\mathbb Z/\ell}(G_{\bar\eta})
-\sum_{s\in S}\varepsilon_s(G),
\]
donde, para cada $s\in S$, identificando $\bar\eta$ con la clausura
algebraica del cuerpo de fracciones del henselizado de $S$ en $s$, se
tiene
\[
\varepsilon_s(G)=
\dim(G_{\bar\eta})-\dim(G_s)
\sisi{-}{+}\dim\operatorname{Sw}(G_{\bar\eta}).
\]
\footnote{La impresión lleva un signo menos ante el término de Swan.
El signo más viene impuesto tanto por la fórmula estándar del
conductor de Artin citada en la frase precedente como por la identidad
impresa (c), que identifica la suma alternada con la dimensión
\emph{total} $\dim+\dim\operatorname{Sw}$.}

Para $s\in S$, las sucesiones exactas XIII, 2.4.6.3,
\[
\cdots\longrightarrow(R^qf_*\mathbb Z/\ell)_s
\longrightarrow(R^qf_*\mathbb Z/\ell)_{\bar\eta}
\longrightarrow\bigoplus_{x\in F(s)}
\Phi^q_{\bar\eta}(\mathbb Z/\ell)_x
\longrightarrow\cdots
\]
dan
\begin{align*}
&\sum_q(-1)^q\left(
\dim(R^qf_*\mathbb Z/\ell)_{\bar\eta}
-\dim(R^qf_*\mathbb Z/\ell)_s\right)\\
&\hspace{4em}=\sum_{x\in F(s)}\sum_q(-1)^q
\dim\Phi^q_{\bar\eta}(\mathbb Z/\ell)_x,
\end{align*}
y
\[
\sum_q(-1)^q\dim\operatorname{Sw}
(R^qf_*\mathbb Z/\ell)_{\bar\eta}
=\sum_{x\in F(s)}\sum_q(-1)^q
\dim\operatorname{Sw}\Phi^q_{\bar\eta}(\mathbb Z/\ell)_x.
\]
Por consiguiente,
\begin{equation}\tag{c}
\sum_q(-1)^q\varepsilon_s(R^qf_*\mathbb Z/\ell)
=\sum_{x\in F(s)}\dim\operatorname{tot}\Phi_x.
\end{equation}
La fórmula de Euler--Poincaré se reduce, pues, a
\begin{align*}
\chi(X)
&=\sum_q(-1)^q\chi(S,R^qf_*\mathbb Z/\ell)\\
&=\sum_q(-1)^q\left(
\chi(S)\dim(R^qf_*\mathbb Z/\ell)_{\bar\eta}
-\sum_s\varepsilon_s(R^qf_*\mathbb Z/\ell)\right)\\
&=\chi(S)\sum_q(-1)^q
\dim(R^qf_*\mathbb Z/\ell)_{\bar\eta}
-\sum_s\sum_q(-1)^q\varepsilon_s(R^qf_*\mathbb Z/\ell)\\
&=\chi(S)\chi(X_{\bar\eta})
-\sum_s\sum_{x\in F(s)}\dim\operatorname{tot}\Phi_x\\
&=\chi(S)\chi(X_{\bar\eta})
-(-1)^n\sum_{x\in F}n(X/S,x).
\end{align*}

\pageoriginale

Comparando esta fórmula con (2.1.1), se obtiene 2.2.

\subsubsection*{Corolario 2.3.}
\underline{Bajo las hipótesis y con las notaciones precedentes,
supongamos que $x_0\in F$ y que, para todo $x\in F$ distinto de
$x_0$, la fibra de $f$ que pasa por $x$ presenta en $x$ una
singularidad cuadrática ordinaria. Entonces}
\[
\mu(X/S,x_0)=n(X/S,x_0).
\]
Para $x\ne x_0$, la Proposición 1.13 da
\[
\mu(X/S,x)=n(X/S,x),
\]
y basta aplicar 2.2.

\subsubsection*{Teorema 2.4.}
\underline{Sean $S$ un trait henseliano de igual característica con
cuerpo residual perfecto, $f:X\to S$ un morfismo plano de tipo finito
y $\bar x$ un punto geométrico de $X$ cuya imagen en $X$ es un punto
cerrado $x$ de $X_s$ que es un punto aislado de no lisitud de $X/S$.
Supongamos que $X$ es regular en $x$. Entonces el número total (1.8)
de ciclos evanescentes de $X/S$ en $\bar x$ coincide con el número de
Milnor $\mu(X/S,x)$.}

Sea $S'$ el trait henseliano deducido de $S$ extendiendo el cuerpo
residual $k(s)$ a su clausura separable, y pongamos
$X'=X\times_{\sisi{s}{S}}S'$. La afirmación 2.4 para $X/S$ equivale
entonces a la afirmación para $X'/S'$. Como $k(s)$ ya se supone
perfecto, podemos suponer que el cuerpo residual de $S$ es
algebraicamente cerrado.

Para tal $S$, sea $S'$ su completado y pongamos
$X'=X\times_{\sisi{s}{S}}S'$. El número de Milnor en $x$ es el mismo
para $X/S$ y $X'/S'$, y por XIII, 2.4.2 lo mismo vale para el número
total de ciclos evanescentes. Podemos suponer además que $S$ es
completo, por tanto isomorfo a $k[[t]]$, con $k$ algebraicamente
cerrado. Si $X_s$ tiene dimensión $n$ en $x$, se tiene:

\subsubsection*{Proposición 2.5.}
\underline{Existen una variedad proyectiva lisa $Y$ sobre $k$,
puramente de dimensión $n+1$, un morfismo
$f:Y\to\mathbb P^1_k$ liso fuera de un conjunto finito $F$}

\pageoriginale

\underline{de puntos de $Y$, y un punto $y_0\in F$, tales que:}
\begin{enumerate}
\item[(a)] \underline{si $y\in F$ y $y\ne y_0$, entonces $y$ es un
punto singular cuadrático ordinario de la fibra de $f$ que pasa por
$y$;}
\item[(b)] \underline{$f(y_0)=0$ y, si
$a:S\to\mathbb P^1_k$ es el isomorfismo natural entre $S$ y el
completado de $\mathbb P^1_k$ en $0$, existe un $S$-isomorfismo entre
el henselizado de $X$ en $x$ y el henselizado de $a^*Y$ en $y_0$.}
\end{enumerate}

Mostremos que $2.5\Rightarrow2.4$. Por 2.3 y 2.5(a), el número de
Milnor de $\sisi{Y/T}{Y/\mathbb P^1_k}$ en $y_0$ coincide con el
número total de ciclos evanescentes. Como antes, XIII, 2.4.2 da el
mismo resultado para $a^*Y/S$, y por tanto también para $(X/S,x)$,
que por 2.5(b) es localmente isomorfo a $(a^*Y/S,y_0)$ para la
topología étale.

Demostremos 2.5. El completado del anillo local $\mathcal O_{X,x}$ es
isomorfo, como $k[[t]]$-álgebra, a un cociente
\[
\widehat{\mathcal O}_{X,x}
=k[[t,x_1,\ldots,x_{n+1}]]/(f).
\]
El teorema de las funciones implícitas XV, 1.1.2 (es decir, [1],
Lema 5.10) muestra entonces que existe un entero $N$ tal que la
condición 2.5(b) se cumple siempre que
\[
\widehat{\mathcal O}_{Y,y_0}
\simeq k[[t,x_1,\ldots,x_{n+1}]]/(f_1)
\]
como $k[[t]]$-álgebras, con $f_1-f$ en la potencia $N$-ésima
$\mathfrak m^N$ del ideal maximal de
$k[[t,x_1,\ldots,x_{n+1}]]$.

Sean $(t_0,t_1)$ y $(x_0,\ldots,x_{n+1})$ los sistemas usuales de
coordenadas homogéneas sobre $\mathbb P^1$ y $\mathbb P^{n+1}$.
Tomaremos por $Y/\mathbb P^1$ la subvariedad de
$\mathbb P^1\times\mathbb P^{n+1}$ definida por una ecuación
\begin{equation}\tag{2.5.1}
g(t_0,t_1;x_0,\ldots,x_{n+1})=0,
\end{equation}
homogénea de grado $d'$ en $(t_0,t_1)$ y de grado $d''$ en los
$x_i$. En términos más intrínsecos, $g$ será una sección de
\[
\mathcal O(d',d'')=\operatorname{pr}_1^*\mathcal O(d')
\otimes\operatorname{pr}_2^*\mathcal O(d'').
\]

\pageoriginale

La condición (b) se cumplirá si
\begin{equation}\tag{2.5.2}
g(1,t;1,x_1,\ldots,x_{n+1})\equiv f\pmod{\mathfrak m^N}.
\end{equation}

Fijemos $d'$ y $d''$. Los $g$ del tipo (2.5.1) son entonces los puntos
racionales de un espacio afín $A$ sobre $k$. Para
$u\in\mathbb P^1\times\mathbb P^{n+1}$, designemos por
$B(u,\sisi{k}{r})$ el espacio afín cuyos puntos son los
$\sisi{k}{r}$-jets de secciones de $\mathcal O(d',d'')$ en $u$.
Si $u_0=(1,0;1,0,\ldots,0)$ y $u\ne u_0$, entonces, para $d'$ y
$d''$ suficientemente grandes, la aplicación de restricción
\begin{equation}\tag{2.5.3}
A\longrightarrow B(u_0,N)\times B(u,2),
\end{equation}
siempre lineal, es sobreyectiva.

Sea $C(u)\subset B(u,2)$ el conjunto de los $2$-jets que son el
$2$-jet del germen cero o de un germen $\widetilde h$ tal que
\begin{enumerate}
\item[(a)] $u$ pertenece a la variedad $V$ definida por
$\widetilde h=0$;
\item[(b)] $u$ no es ni un punto liso ni un punto cuadrático ordinario
de la fibra de $V/\mathbb P^1$ que pasa por $u$.
\end{enumerate}
Se comprueba fácilmente que $C(u)$ es cerrado de codimensión $n+3$
en $B(u,2)$.

Sea $D$ el subespacio afín de $A$ formado por los $g$ que satisfacen
(2.5.2). Para $u\ne u_0$, sea $D(u)$ el subespacio de $D$ formado por
el cero y por los $g$ tales que $g(u)=0$ y $u$ no es ni un punto liso
ni un punto cuadrático ordinario de la fibra de la variedad $V$
definida por $g=0$ que pasa por $u$. Si $d'$ y $d''$ son
suficientemente grandes para que (2.5.3) sea sobreyectiva, entonces
$D(u)$ tiene codimensión $n+3$ en $D$. La unión de los $D(u)$ para
$u\ne u_0$ tiene, por tanto, codimensión al menos
\[
1=(n+3)-(n+2),
\]
y un elemento general $g$ de $D$ no pertenece a ninguno de los
$D(u)$. Tal elemento resuelve el problema.

\pageoriginale

\subsection*{Bibliografía}

\begin{enumerate}[label={[\arabic*]},leftmargin=3em]
\item M. Artin, ``Algebraic approximation of structures over complete
local rings'', \emph{Publ. Math. IHES} \underline{36} (1969)
(volumen de Zariski).
\item Lê Dũng Tráng, \emph{Singularités isolées des hypersurfaces
complexes}, Publ. Centre Math. École Polytechnique, mayo de 1969.
\item J. Milnor, \emph{Singular Points of Complex Hypersurfaces},
Princeton University Press, 1968.
\item M. Raynaud, ``Caractéristique d'Euler--Poincaré d'un faisceau et
cohomologie des variétés abéliennes'', Séminaire Bourbaki 286,
febrero de 1965.
\item J.-P. Serre, \emph{Corps locaux}, Publ. Inst. Math. Univ.
Nancago VIII, Hermann, 1962.
\item J.-P. Serre, ``SGA 5, IX. Introduction à la théorie de Brauer'',
reproducido en \emph{Représentations linéaires des groupes},
3.\textsuperscript{a} parte, Hermann, París, 1967.
\end{enumerate}
